\documentclass{article}
\title{Lightweight Anomaly Detection for Edge Pump-Station Sensor Streams}
\author{Synthetic Authors (demo data)}
\begin{document}
\maketitle

\begin{abstract}
Pump stations produce continuous sensor streams that must be screened for faults close to the equipment. We present a lightweight detector based on an isolation forest that runs on low-power edge hardware. The detector reaches an F1 of 0.91 on our pump-station data while using little memory.
\end{abstract}

\section{Introduction}
Water utilities operate hundreds of pump stations and need early warning of bearing wear, blockages and sensor faults. Sending every reading to a central server is costly, so detection should happen on the edge.

\section{Related Work}
Classical statistical process control flags readings that leave fixed control limits \cite{shewhart}. Learning-based detectors such as isolation forests \cite{liu2008} isolate anomalies with random partitions and need no labels.

\section{Method}
The detector scores a sliding window of recent readings with an isolation forest. A reading is reported as anomalous when its score exceeds a threshold that is tuned on a validation split. The model is small enough to run on a microcontroller-class device.

\section{Experiments}
\subsection{Setup}
We evaluate the detector on sensor data from pump stations. We compare against standard baselines. All results come from a single run on the test split.

\subsection{Results}
Our detector reaches an F1 of 0.91 on the test split. Precision and recall are balanced, and the detector raised few false alarms during normal operation.

\section{Conclusion}
We presented a lightweight isolation forest detector for edge pump-station monitoring and showed that it detects faults accurately on our data.

\begin{thebibliography}{9}
\bibitem{shewhart} W. Shewhart. Economic control of quality of manufactured product. 1931.
\bibitem{liu2008} F. Liu, K. Ting, Z. Zhou. Isolation forest. ICDM 2008.
\end{thebibliography}
\end{document}
